\documentclass[11pt]{article}

\usepackage[T1]{fontenc}
\usepackage[utf8]{inputenc}
\usepackage[spanish,es-nodecimaldot]{babel}
\usepackage{lmodern}
\usepackage[margin=1in]{geometry}
\usepackage{amsmath,amssymb}

\usepackage{hyperref}
\hypersetup{
  colorlinks=true,
  urlcolor=blue,
  linkcolor=black,
  citecolor=black
}

\setlength{\parindent}{0pt}
\setlength{\parskip}{0.45em}
\renewcommand{\labelenumi}{\textnormal{(\alph{enumi})}}
\pagestyle{plain}

\let\enumerateoriginal\enumerate
\let\endenumerateoriginal\endenumerate
\renewenvironment{enumerate}{%
  \enumerateoriginal
  \setlength{\itemsep}{0.5em}%
  \setlength{\parsep}{0pt}%
  \setlength{\topsep}{0.4em}%
}{\endenumerateoriginal}

\newcounter{problema}
\newenvironment{problema}[1]{%
  \par
  \refstepcounter{problema}%
  \addvspace{0.9em}%
  \noindent\textbf{Ejercicio \theproblema. #1}
  \par\nobreak\smallskip
}{\par}

\newcommand{\R}{\mathbb{R}}
\newcommand{\MRS}{\mathrm{MRS}}
\newcommand{\MPN}{\mathrm{MPN}}

\begin{document}

\begin{center}
  {\large\textbf{Centro de Investigación y Docencia Económicas}}\\[0.25em]
  {\large\textbf{Macroeconomía II}}\\[0.2em]
  {\large\textbf{Problem Set 2: Equilibrio General}}\\[0.35em]
  LECO -- Otoño 2026
\end{center}

\vspace{0.6em}
\hrule
\vspace{0.6em}

\begin{problema}{Equilibrio general en una economía cerrada de un periodo}
Considere una economía con un hogar representativo, una empresa representativa y un gobierno. El hogar y la empresa se comportan de forma racional, en un entorno de competencia perfecta. El bien de consumo es el bien numerario.

El hogar dispone de una unidad de tiempo, que distribuye entre trabajo y ocio. Elige consumo $C>0$ y ocio $\ell\in(0,1)$, y ofrece $N^s=1-\ell$ unidades de trabajo. Sus preferencias están representadas por la siguiente función de utilidad
\[
  U(C,\ell)=\log C+\psi\log\ell,
  \qquad \psi>0.
\]

La empresa demanda $N^d$ unidades de trabajo y produce de acuerdo con la tecnología
\[
  Y=z\sqrt{N^d},
  \qquad z>0.
\]
El hogar es propietario de la empresa y recibe todos sus beneficios.

El gobierno compra $G$ unidades del bien de consumo y financia este gasto mediante un impuesto de suma fija $T>0$ cobrado al hogar. El presupuesto público está balanceado: $ G=T.$ 

Para que exista una solución interior, suponga que $0< G<z$.

\begin{enumerate}
  \item Defina formalmente un equilibrio competitivo en esta economía.

  \item Obtenga las funciones de oferta laboral y demanda de consumo del hogar, así como las funciones de demanda laboral, oferta de bienes y beneficios de la empresa. 

  \item Encuentre el equilibrio competitivo. Esto es, la asignación $(N^*,\ell^*,C^*,Y^*, \pi^*, T^*)$ y el vector de precios $(1,w^*)$.

  \item Muestre que, una vez satisfechas la restricción presupuestaria del hogar, la identidad de beneficios de la empresa, el presupuesto del gobierno y el vaciado del mercado de trabajo, el mercado de bienes también se vacía: $Y^*=C^*+G$.

  \item Formule el problema del planificador social para esta economía y obtenga sus condiciones de primer orden. Explique por qué la asignación competitiva coincide con la asignación del planificador.

  \item Grafique el equilibrio competitivo en los siguientes tres diagramas:
  \begin{enumerate}
    \renewcommand{\labelenumii}{\textnormal{(\roman{enumii})}}
    \item mercado de trabajo, en el plano $(N,w)$;
    \item mercado de bienes, representando la producción $Y^s(w)$ y la demanda agregada $C^d(w)+G$ como funciones del salario real;
    \item el problema de asignación en el plano $(\ell,C)$, mostrando la frontera de posibilidades
    de producción y la curva de indiferencia que pasa por el equilibrio.
  \end{enumerate}
  Explique por qué el segundo diagrama no representa un mecanismo independiente de ajuste de precios.

  \item \textbf{Choque de gasto público.}
 Ante el deterioro del entorno macroeconómico internacional provocado por la crisis financiera global, el Gobierno de México anunció en octubre de 2008 el \href{https://www.eleconomista.com.mx/el-empresario/Felipe-Calderon-lanza-plan-anticrisis-financiera-20081009-0003.html}{Programa para Impulsar el Crecimiento y el Empleo (PICE)}, el cual se incorporó al Paquete Económico 2009. Entre sus medidas se encontraba un programa de \href{https://expansion.mx/expansion/2008/12/02/cartera-abierta}{gasto público, principalmente en infraestructura, por 65.1 mil millones de pesos}. La SHCP informó que el estímulo fiscal del programa equivaldría al 0.7 por ciento del PIB.\footnote{Véanse Secretaría de Hacienda y Crédito Público, \emph{Comunicado de prensa 079/2008}, 8 de octubre de 2008, e \emph{Informe de Labores de la SHCP 2009}. Ambos documentos están disponibles en el portal institucional de la SHCP.}
  \begin{enumerate}
    \renewcommand{\labelenumii}{\textnormal{(\roman{enumii})}}
    \item De acuerdo con el modelo, ¿cuál es el impacto de este choque sobre la economía? Represente gráficamente la respuesta de las variables endógenas, los mecanismos de transmisión e interprete económicamente.
    \item Si tuviera que evaluar la predicción del modelo con datos de la economía mexicana, ¿qué dirección esperaría observar en el PIB, el consumo privado, el empleo y el salario real a partir del estímulo fiscal? 
  \end{enumerate}

  \item \textbf{Choque de productividad.} Suponga que la economía experimenta un aumento de productividad, de $z$ a $z'>z$.
  \begin{enumerate}
    \renewcommand{\labelenumii}{\textnormal{(\roman{enumii})}}
    \item De acuerdo con el modelo, ¿cuál es el impacto de este choque sobre la economía? Represente gráficamente la respuesta de las variables endógenas, los mecanismos de transmisión e interprete económicamente.
    \item Compare los comovimientos producidos por el choque de productividad con los generados por el choque fiscal. ¿Qué variables permiten distinguir ambos choques dentro del modelo?
  \end{enumerate}

  \item \textbf{Identificación empírica de los choques.}
  Construya una tabla con los signos de respuesta de $C$, $Y$, $N$, $\ell$, $w$ y $\pi$ ante los dos choques estudiados. Suponga que, en los datos, se observa un aumento simultáneo de producción, consumo y salario real. ¿Cuál de los dos choques es más consistente con ese patrón? ¿Cambiaría su
  respuesta si la producción aumentara mientras el consumo y el salario real disminuyen? Explique qué información adicional necesitaría para interpretar estos comovimientos como evidencia de alguno de los choques.
  \end{enumerate}

  \end{problema}

\end{document}
